%% test-tagging.tex
%% ================
%% A test file run with \DocumentMetadata{tagging=on}: the footnote
%% texts are typeset where their marks are and output after the box, so
%% that in the structure tree each footnote (FENote) directly follows
%% its mark (Lbl), as for \footnote in the body text.
%%
%%   pdflatex test-tagging.tex   (or xelatex / lualatex)
%%   platex  -jobname=test-tagging "\def\CLS{jarticle}\input test-tagging"
%%   (then dvipdfmx test-tagging)
%%   texlua test-tagging.lua test-tagging.pdf
%% The file to run is test-basic.tex by default; set \TPFfile for another
%% one (test-breakable, test-japanese, ...), and \TPFhyperref for
%% test-basic.tex with hyperref.
%%
%% What to check:
%%   * the same as for the file that is run;
%%   * test-tagging.lua: every footnote mark is paired with its FENote,
%%     except the two \footnotemark of test-basic.tex whose texts are
%%     given later by \footnotetext.
%% pLaTeX and upLaTeX do not write a structure tree; only the first
%% point is checked with them.
\DocumentMetadata{tagging=on}
\ifdefined\TPFfile\else\def\TPFfile{test-basic}\fi
\input{\TPFfile}
